\documentclass[11pt]{amsart}
\usepackage{amsmath,amssymb,amsthm,booktabs}
\usepackage[margin=1.13in]{geometry}
\usepackage[hidelinks]{hyperref}
\DeclareMathOperator{\ord}{ord}
\DeclareMathOperator{\tr}{tr}
\newcommand{\Z}{\mathbb Z}
\newcommand{\eps}{\varepsilon}
\newcommand{\sq}{\mathbin{\square}}
\newcommand{\NB}{\operatorname{nb}}
\newtheorem{theorem}{Theorem}[section]
\newtheorem{lemma}[theorem]{Lemma}
\newtheorem{corollary}[theorem]{Corollary}
\newtheorem{proposition}[theorem]{Proposition}
\theoremstyle{definition}
\newtheorem{definition}[theorem]{Definition}
\theoremstyle{remark}
\newtheorem{remark}[theorem]{Remark}
\newtheorem{question}[theorem]{Question}
\title[Nonbacktracking and primitive cycle symbols]{Nonbacktracking and primitive cycle symbols\\ on the Cartesian graph ring}
\date{Draft for mathematical review --- October 23, 2026}
\begin{document}
\begin{abstract}
For finite simple graphs under disjoint union and Cartesian product, we
study nonbacktracking Hashimoto traces, oriented primitive cyclic classes,
and ordinary non-induced simple cycles using centered augmentation-ideal
symbols. A coordinate-support argument proves, at every $k\ge2$, that the
nonbacktracking trace of length $2k$ has differential order exactly $k$
and pure top symbol $\beta_k d^{\otimes k}$, with
$\beta_k=2^k k!R_k>0$ and $R_k$ a finite count of pairings of cyclic
positions avoiding adjacent pairs. M\"obius inversion yields the primitive
symbol $\beta_k/(2k)d^{\otimes k}$. An independent coordinate-support
classification gives the simple-cycle top symbol
$\gamma_k d^{\otimes k}$, where $\gamma_k$ counts simple full-support
cycles in a $k$-cube and is strictly positive. Thus the aggregate
repeated-vertex deficit also has a pure top symbol. In particular, the
correct value is $\gamma_7=880\,865\,280$, not the negative coefficient
previously reported from an incomplete enumeration of $Q_7$.
\end{abstract}
\maketitle
\tableofcontents

\section{Setting and proof status}
Let $\mathcal A$ be the Grothendieck ring of finite undirected, loopless,
simple graphs, with disjoint union as addition, Cartesian product $\sq$
as multiplication and $K_1$ as unit. Set $\eps(G)=|V(G)|$ and
$d(G)=|E(G)|$, extended additively. Then $d$ is an $\eps$-derivation:
\begin{equation}\label{eq:der}
d(G\sq H)=d(G)\eps(H)+\eps(G)d(H).
\end{equation}
Write $I=\ker\eps$, $u_G=G-\eps(G)K_1$, and, for any additive invariant $f$,
\begin{equation}\label{eq:center}
D_q f(G_1,\ldots,G_q)=f(u_{G_1}\cdots u_{G_q}).
\end{equation}
Our order convention is $\ord_\eps f\le r$ exactly when
$f(I^{r+1})=0$. At exact order $r$, the principal symbol is
$\sigma_r(f)=D_r f$ without a factorial division. This \emph{centered}
convention differs from the uncentered iterated Leibniz defect used in
some historical drafts when $f(K_1)\ne0$.

The generic Cartesian semiring, character and differential filtration are
standard constructions; adjacency Kronecker sums and Bass--Ihara theory
are imported ingredients, not claims of novelty. The proof statements below
are a research manuscript under mathematical and prior-art review, with
repository evidence at \cite{GPR}. The historical Paper III v4 is kept as
an archival source rather than silently corrected in place.

\section{Hashimoto traces and centered coordinate support}
For a graph $G$ let $\vec E(G)$ be its directed arcs. The Hashimoto
matrix $B_G$ has entry $B_{(u,v),(v,w)}=1$ if $w\ne u$ and zero otherwise.
Define
\begin{equation}\label{eq:nb}
b_m(G):=\tr(B_G^m),\qquad m\ge1.
\end{equation}
This is the count of \emph{rooted oriented cyclically nonbacktracking}
closed walks of length $m$: immediate reversal is excluded, including at
the wraparound between the last and first edge. It is additive under
disjoint union and $b_m(K_1)=0$.

\begin{lemma}[Centered support defect]\label{lem:support}
For actual graphs $G_1,\ldots,G_q$, $D_q b_m(G_1,\ldots,G_q)$ is exactly the
number of rooted oriented cyclically nonbacktracking closed length-$m$
walks in $G_1\sq\cdots\sq G_q$ whose projections change every coordinate.
\end{lemma}
\begin{proof}
Every product edge changes exactly one coordinate. A walk has a unique set
$S\subseteq[q]$ of coordinates used. A term indexed by
$T\subseteq[q]$ in the expansion of~\eqref{eq:center} counts walks whose
used-coordinate set is contained in $T$, with arbitrary fixed coordinates
outside $T$. The cyclic nonbacktracking condition is invariant under this
fiber inclusion. The coefficient of a particular walk with used set $S$ is
$\sum_{T\supseteq S}(-1)^{q-|T|}$, equal to $1$ when $S=[q]$ and $0$
otherwise. The empty-index term vanishes because $b_m(K_1)=0$.
No vertex-injectivity assumption was used.
\end{proof}

\begin{lemma}[Uniform filtration bound]\label{lem:bound}
For every $m\ge1$, $b_m(I^{\lfloor m/2\rfloor+1})=0$.
\end{lemma}
\begin{proof}
The projection of a closed walk onto any coordinate it uses is a
nonconstant closed walk in a loopless graph and traverses at least two
edges (possibly one edge twice). The disjoint coordinate-support sets of
steps therefore satisfy $m\ge 2|S|$. By Lemma~\ref{lem:support} all
$D_qb_m$ with $q>\lfloor m/2\rfloor$ vanish on actual factor graphs.
The $u_G$ additively generate $I$; multilinearity extends the identity
to all of $I^{\lfloor m/2\rfloor+1}$.
\end{proof}

\section{A pure all-order nonbacktracking symbol}
Let $W_k$ be the set of words of length $2k$ over
$\{1,\ldots,k\}$ in which each letter occurs twice and consecutive
letters, \emph{including the cyclic last--first pair}, differ. Put
\begin{equation}\label{eq:beta}
R_k=\frac{|W_k|}{k!},\qquad \beta_k=2^k|W_k|=2^k k!R_k.
\end{equation}
The integer $R_k$ counts pair partitions of $2k$ positions in which no
pair joins cyclically adjacent positions.

\begin{theorem}[Global nonbacktracking symbol]\label{thm:nb}
On the \emph{full} graph ring $\mathcal A$, for every $k\ge2$,
\begin{equation}\label{eq:nb-symbol}
\ord_\eps(b_{2k})=k,\qquad
\sigma_k(b_{2k})=\beta_k d^{\otimes k},\qquad \beta_k>0.
\end{equation}
\end{theorem}
\begin{proof}
The order bound follows from Lemma~\ref{lem:bound}. In a closed walk
of length $2k$ changing all $k$ product coordinates, each coordinate is
changed exactly twice. Its two-step projection crosses one uniquely
determined undirected edge and returns on the same edge. Consequently
every full-support walk determines exactly one edge in each factor and
is confined to their Cartesian product, a copy of $Q_k$.

Conversely, choose an edge from every $G_i$, choose the starting endpoint
of each selected edge ($2^k$ choices), and choose a word in $W_k$.
Flipping the selected coordinates according to the word produces a closed
walk, as each coordinate occurs twice. A pair of successive steps
backtracks if and only if their coordinate labels agree; this holds also
at the cyclic seam. Thus these choices parametrize exactly the
full-support cyclically nonbacktracking rooted walks, even when vertices
are revisited nonconsecutively. The selected edge tuple is recovered
uniquely from the walk. Lemma~\ref{lem:support} gives
\[
D_kb_{2k}(G_1,\ldots,G_k)
=2^k |W_k|\prod_i d(G_i)=\beta_k\prod_i d(G_i).
\]
For every $k\ge2$, the word
$1,2,\ldots,k,1,2,\ldots,k$ lies in $W_k$, including the seam.
Thus $\beta_k>0$ and evaluation on $k$ copies of $K_2$ proves
$\ord_\eps(b_{2k})\ge k$.
\end{proof}

\begin{proposition}[Finite coefficient formula]\label{prop:finite}
For $k\ge2$, with $(-1)!!=1$,
\begin{equation}\label{eq:inclusion}
R_k=\sum_{j=0}^k(-1)^j
\frac{2k}{2k-j}\binom{2k-j}{j}(2k-2j-1)!!.
\end{equation}
This equals the historical Chebyshev--moment leading coefficient
$a_k^{\rm nb}$, so its nonvanishing is a purely finite statement.
\end{proposition}
\begin{proof}
The forbidden pairs are the $2k$ edges of the cycle $C_{2k}$
formed by consecutive position labels. Inclusion--exclusion over a
$j$-subset of forbidden pairs is nonzero only when those pairs are
disjoint. The number of $j$-matchings of $C_{2k}$ is
$\frac{2k}{2k-j}\binom{2k-j}{j}$. Once the $j$ pairs are fixed, the
remaining $2k-2j$ vertices can be matched in
$(2k-2j-1)!!$ ways. This proves~\eqref{eq:inclusion}.
The coefficients of the regular-graph Chebyshev polynomial
$\phi_{2k}(\lambda;q)$ are
$(-1)^j\frac{2k}{2k-j}\binom{2k-j}{j}q^j\lambda^{2k-2j}$;
substituting Gaussian even \emph{moments} $(2k-2j-1)!!$ gives
exactly~\eqref{eq:inclusion}. This identifies the scalar formula, not
an additional assumption about the global symbol.
\end{proof}

\section{Primitive cyclic classes}
Let $p_m(G)$ count oriented primitive cyclic equivalence classes of
cyclically nonbacktracking length-$m$ walks, modulo cyclic rotation but
\emph{not} orientation reversal. A trace term is an indexed arc sequence, equivalently a vertex traversal
with specified initial time and direction; there is no extra rooting
factor in~\eqref{eq:nb}. A primitive class of length $r$ has exactly $r$
distinct rooted sequences, even when raised to a power. Periodicity must
be tested on the actual arc or vertex sequence, not merely its coordinate
word (the square word $1212$ is a primitive traversal).
Powers of primitive classes give the
standard identity
\begin{equation}\label{eq:mob}
b_m=\sum_{r\mid m}r p_r,
\qquad p_m=\frac1m\sum_{r\mid m}\mu(m/r)b_r.
\end{equation}
\begin{corollary}[Primitive symbol]\label{cor:primitive}
On $\mathcal A$, for all $k\ge2$,
\[
\ord_\eps(p_{2k})=k,\qquad
\sigma_k(p_{2k})=\frac{\beta_k}{2k}d^{\otimes k}.
\]
\end{corollary}
\begin{proof}
For every proper divisor $r<2k$, Lemma~\ref{lem:bound}
gives $\ord_\eps(b_r)\le\lfloor r/2\rfloor<k$.
Thus the $k$-th centered symbol of~\eqref{eq:mob} has only its
$r=2k$ term. The same identity establishes the upper order bound;
Theorem~\ref{thm:nb} proves nonvanishing.
\end{proof}

\section{Ordinary simple cycles and the aggregate deficit}
Write $c_{2k}(G)=\mathrm{inj}(C_{2k},G)/(4k)$ for the number of
ordinary, not necessarily induced, simple cycle copies. Here
\emph{simple} means that all vertices of a traversal are distinct before
the final return. Let $A_k$ be the number of pair partitions of
$\{1,\ldots,2k\}$ having no nonempty proper contiguous interval
that is a union of whole pairs.

\begin{theorem}[Simple even-cycle symbol]\label{thm:cycles}
On the full graph ring, for every $k\ge2$,
\[
\ord_\eps(c_{2k})=k,\qquad
\sigma_k(c_{2k})=\gamma_k d^{\otimes k},\qquad
\gamma_k=\frac{2^k k! A_k}{4k}>0.
\]
\end{theorem}
\begin{proof}
The centered defect $D_q c_{2k}$ counts ordinary cycle copies using
all $q$ Cartesian factors, by the same used-coordinate inclusion--exclusion
as Lemma~\ref{lem:support}. Every used coordinate must change at least
twice, giving order at most $k$. Maximal-support cycles project to an
edge out and back in each factor, so they are uniquely carried by a
chosen edge tuple and a full-support simple cycle in $Q_k$.

A full-support cube traversal yields a pairing of its two occurrences
of each coordinate. Two visited cube vertices coincide if and only if
the intervening flips consist of complete pairs. Hence vertex
injectivity up to the final return is exactly the condition defining
$A_k$. There are $k!$ coordinate relabelings and $2^k$ starts; each
unoriented unrooted $2k$-cycle is counted $4k$ times. This gives the
coefficient. The word $1,2,\ldots,k,1,2,\ldots,k$ is a simple
full-support cycle, hence $\gamma_k>0$ and the order is exact.
\end{proof}

Define the oriented non-simple primitive deficit
\[
\Delta_k=p_{2k}-2c_{2k}.
\]
\begin{corollary}[Aggregate top correction]\label{cor:deficit}
For every $k\ge2$,
\begin{equation}\label{eq:deficit}
D_k\Delta_k=\eta_k d^{\otimes k},\qquad
\eta_k=\frac{2^k k!}{2k}(R_k-A_k)\ge0.
\end{equation}
\end{corollary}
\begin{proof}
Here $D_k\Delta_k$ denotes the degree-$k$ centered defect, which may
vanish; we do not assert that $\Delta_k$ has exact order $k$.
Subtract twice the simple-cycle top symbol in
Theorem~\ref{thm:cycles} from the primitive symbol in
Corollary~\ref{cor:primitive}. Every admissible simple-cycle
word is cyclically nonbacktracking, so $A_k\le R_k$.
\end{proof}
\begin{remark}
The full graph ring distinction matters: $\Delta_3$ need not vanish
as a \emph{function} on all graphs (two triangles may share a vertex),
even though its order-$3$ top coefficient is zero. On the bipartite
sector $\Delta_k$ vanishes identically for $k\le3$.
No assertion here classifies non-simple primitive cycles by a partition
of their constituent simple-cycle lengths; the lower-order
junction multiplicities remain unresolved.
\end{remark}

\section{Exact coefficient ledger and correction to the historical draft}
\begin{center}
\begin{tabular}{@{}r r r r r r@{}}\toprule
$k$ & $R_k$ & $\beta_k$ & $A_k$ & $\gamma_k$ & $\eta_k$\\\midrule
2&1&8&1&1&0\\
3&4&192&4&16&0\\
4&31&11\,904&27&648&192\\
5&293&1\,125\,120&248&47\,616&17\,280\\
6&3\,326&153\,262\,080&2\,830&5\,433\,600&1\,904\,640\\
7&44\,189&28\,507\,207\,680&38\,232&880\,865\,280&274\,498\,560\\
\bottomrule\end{tabular}
\end{center}
The values of $R_k$ can be recomputed from the exact finite
inclusion--exclusion in~\eqref{eq:inclusion}. The $A_k$ counts and
ordinary-cycle cube enumeration are independently checked in the
repository, with separate coordinate-word and edge-subgraph oracles
\cite{GPR}. The prior draft claimed
\[
 c_{14}(Q_7)=406\,517\,760,\qquad \gamma_7=-2\,049\,830\,400.
\]
Those values are false under ordinary simple-cycle conventions. The
corrected finite enumeration is
\begin{equation}\label{eq:correct7}
\boxed{c_{14}(Q_7)=3\,337\,213\,440,\qquad
\gamma_7=880\,865\,280>0.}
\end{equation}
Indeed, writing $B_j$ for full-$j$-coordinate $14$-cycles in $Q_j$,
the unique minimal-face decomposition gives
\[
 c_{14}(Q_7)=\sum_{j=0}^7\binom7j2^{7-j} B_j,
\quad B_7=\gamma_7=880\,865\,280.
\]
The previous value omitted full-support cycles and generated a
spurious negative seventh finite difference. Theorem~\ref{thm:cycles}
proves $\gamma_k>0$ for \emph{every} $k\ge2$, so the claimed
sign reversal and possible sign oscillation are withdrawn.

\section{Verification scope and remaining questions}
The checked-in test suite compares independent directed-arc traces
against centered defects on hypercubes and on products involving
nonregular and nonbipartite graphs, and checks cyclic-seam conditions,
M\"obius factors and the finite matching formula. A separate finite
replay recomputed
\[
 (R_2,\ldots,R_7)=(1,4,31,293,3326,44189),
\]
and, for example, the centered two-factor nonbacktracking symbol on
$(P_3,K_3)$ at length $4$ equals $48=8\cdot2\cdot3$; the
three-factor symbol on $(P_3,K_2,C_4)$ at length $6$ is
$1536=192\cdot2\cdot1\cdot4$.
These tests corroborate counting conventions; the all-order conclusions
rest on the explicit bijections and augmentation-filtration proofs.

\begin{question}
Can the aggregate correction in~\eqref{eq:deficit} be refined into a
\emph{proved} partition-indexed junction decomposition, with separately
identified $\mu_\lambda$ and a transfer object for repeated-vertex
incidence? The coefficients at $k=6$ alone do not determine its three
suggested partition types.
\end{question}
\begin{question}
Which of the operator/symbol formulations survive a claim-by-claim
prior-art comparison with the Ihara--Bass and Cartesian graph-bundle
literature? No originality conclusion follows from this manuscript.
\end{question}

\begin{thebibliography}{9}
\bibitem{FoataZeilberger}
D. Foata and D. Zeilberger, \emph{A combinatorial proof of Bass's
evaluations of the Ihara--Selberg zeta function for graphs},
Trans. Amer. Math. Soc. \textbf{351} (1999), 2257--2274.
\bibitem{RG}
\emph{A ring-theoretic framework for graph invariants via Cartesian
products}, parent manuscript (historical research draft, 2026).
\bibitem{GPR}
Graph-product-ring-research, \emph{Ordinary even cycles: support,
augmented order, and principal symbol}; \emph{Nonbacktracking/primitive
symbols by coordinate support}; and exact-cube verification oracles,
2026. Repository research records, pending independent acceptance.
\end{thebibliography}
\end{document}
